\section{相关工作}
\label{sec:related}
\paragraph{NBV 与主动场景采集。}
Bircher 等利用滚动时域 NBV 规划，根据新可观测空间评价探索分支\cite{bircher2016nbv}。
Surface Edge Explorer 则根据已有三维观测的密度与边界引导后续表面采集\cite{border2018see}。
这些工作使用不同表示回答追加测量的位置选择问题。
本文并不认为通用探索目标不适用于建图，而是研究当终点变为另一台机器人的操作任务时，
其空间加权是否仍然合适。因此，本文的 Generic 对照是在同一规划器中的目标函数级控制，
并非对上述完整已发表系统的复现或直接性能对比。

\paragraph{任务导向感知与主动抓取。}
Gualtieri 与 Platt 研究视点选择对抓取定位与置信度的影响\cite{gualtieri2017viewpoint}。
Breyer 等将遮挡目标区域的探索、持续更新的抓取预测，以及执行或继续观测的决策相结合
\cite{breyer2022target}。这些工作说明测量价值可以取决于抓取目标，而不仅是地图完整性。
本文将这一思想连接到异构机器人之间：空中观测由地面移动机械臂的可用站位支撑加权，
并显式考虑底盘完整占地范围。地面栅格的相关性因此来自未来站位，
而不仅来自其是否属于未完整观测的目标。本文不主张首次提出任务导向感知。

\paragraph{空地协作。}
ColAG 将空中感知与建图连接到感知受限地面车辆的导航需求，并包含空中航点调度
\cite{li2023colag}。Lenz 等展示了面向砌墙的空地系统自主感知与砖块操作
\cite{lenz2020wall}。这些工作说明协同感知与物理动作需要整体考虑。
本文侧重于连接地面机械臂操作支撑与空中视点评价的中间表示，
并通过物理取回终点与匹配的通用目标比较。
本文不解决一般多机器人任务分配，也不替代协同导航与定位。

\paragraph{可达性与移动操作。}
能力图描述机器人可达工作空间的结构\cite{zacharias2007capability}。
RM4D 提供紧凑表示，支持可达性与逆可达性查询，包括移动机械臂底盘站位应用
\cite{rudorfer2025rm4d}。本文复用 RM4D 生成候选操作支撑，不将其视作本文贡献。
新增环节是把经验证的底盘位姿支撑投影到使用这些位姿前必须感知的环境区域。
可达性本身仍不充分：底盘占据、实测地面支撑、近场可见性与考虑碰撞的执行在后续分别验证。
共用执行接口使这些边界体现在实验结果中，而不假定有效逆运动学解意味着任务完成。
